% `tex.web` §476, the two checks a macro's parameter text gets as it is read:
%
%     if t=zero_token+9 then
%       begin print_err("You already have nine parameters");
%       help2("I'm going to ignore the # sign you just used,")
%         ("as well as the token that followed it."); error; goto continue;
%       end
%     else begin incr(t);
%       if cur_tok<>t then
%         begin print_err("Parameters must be numbered consecutively");
%         help2("I've inserted the digit you should have used after the #.")
%           ("Type `1' to delete what you did use."); back_error;
%         end;
%       cur_tok:=s;
%       end;
%
% Both *recover*.  A tenth `#` and the token after it are dropped; a wrong
% digit is replaced by the right one, with the offending token put back to be
% read as part of the parameter text — `back_error`, so the context carries
% `<to be read again>`.  We rejected the whole definition and said nothing, so
% `\def\a#1#3{x}` defined nothing and left `{x}` to be typeset.
%
% And §399, which the delimited path had and the undelimited one did not:
%
%     loop@+ begin fast_store_new_token(cur_tok); get_token;
%       if cur_tok=par_token then if long_state<>long_call then
%         @<Report a runaway argument and abort@>;
%
% — a `\par` in the argument of a macro that is not `\long` ends the call at
% *whatever depth* the braces are, which is what stops a missing `}` from
% swallowing the rest of the document.  Only the delimiter exemption is
% depth-0, because that is the only place §392 matches one.
\catcode`\{=1 \catcode`\}=2 \catcode`\#=6
\def\a#1#3{x}
\def\b#1#2#3#4#5#6#7#8#9#1{y}
\def\p#1{[#1]}
\long\def\L#1{[#1]}
\p{one

two}
\L{three

four}
\def\q#1\stop{(#1)}
\q one

two\stop
\end
